\begin{frame}{while-Schleife}
\centering\includegraphics[scale=1.0]{images/while}
\end{frame}

\begin{frame}[fragile]{Die \texttt{while}-Schleife}
	Umgangssprachlich:\\
	Wiederhole den Programmcode solange die Bedingung True ist.
	Gefahr: Falls immer das gleiche (Bedingung wird nicht verändert) geprüft wird $=>$ Schleife ohne Ende

    \begin{lstlisting}
while Bedingung == True:
    # Programmcode

    \end{lstlisting}
    
    Beispiel zähle bis 3
    \begin{lstlisting}
i = 1
while i <= 3:
    print(i)
    i = i+1   
    \end{lstlisting}
   
\end{frame}

\begin{frame}[fragile]{while-Schleife Übung}
\textbf{Aufgabe 1}\\
	Verändere den Code so, dass er mit einer while-Schleife 6-mal den \texttt{String} ''Hello World'' ausgibt.

\pause{}
\begin{exampleblock}{Lösung}
	\begin{lstlisting}		 
i = 1
while i <= 6:
	print("Hello World!")
	i = i+1   		
	\end{lstlisting}
\end{exampleblock}
\end{frame}

\begin{frame}[fragile]{while-Schleife Übung}
\textbf{Aufgabe 2}\\
	Schreibe einen Code, der solange fragt ''Nerv ich dich? '' bis er die Eingabe ''JA!'' erhält.\\
	Danach lass das Programm ''Schade'' ausgeben.
	

\pause{}
\begin{exampleblock}{Lösung}
	\begin{lstlisting}		 
		inp = input("Nerv ich dich? ")
		while inp != "JA!":
			  inp = input("Nerv ich dich? ")
		print("Schade")	
	\end{lstlisting}
\end{exampleblock}
\end{frame}



\begin{frame}{Magische Miesmuschel}

	Programmiere einen einfachen Bot.\\ 
	Er wartet mit (''Du: '') auf eine Nein/Ja-Frage und
        beantwortet sie mit einer zufälligen (random) Antwort (''Bot: '' + Antwort). \\
	Er fragt so lange nach bis er ''Ich will nicht mehr'' als Eingabe erhält. Er verabschiedet sich mit "Bot: Bye"\\
	nötig:
	\begin{itemize}
		\item input
		\item while-Schleife
		\item Bedingung
		\item random Zahl (soll einer bestimmten Antwort zugeordnet werden)
	\end{itemize}		
\end{frame}

\begin{frame}[fragile]{Die \texttt{while}-Schleife}
\begin{exampleblock}{Lösung}
\begin{lstlisting}
from random import *

inp = input("Du: ")
while inp != "Ich will nicht mehr" :
	zahl = randint(0, 1)
	if zahl == 0:
		print("Bot: Nein")
	else:
		print("Bot: Ja")
	inp = input("Du: ")
print("Bot: Bye")
\end{lstlisting}
\end{exampleblock}
\end{frame}


